\documentclass[10pt,a4paper]{amsart}
\usepackage[margin=19mm]{geometry}
\usepackage{amsmath,amssymb,mathtools}
\usepackage[scr=rsfs,cal=euler]{mathalfa}
\usepackage{xfrac}
\usepackage[unicode,hidelinks,bookmarks=false]{hyperref}
\newcommand{\ie}{i.e.\ }
\newcommand{\cf}{cf.\ }
\newcommand{\resp}{respectively\ }
\newcommand{\Subsection}{\subsection}
\providecommand{\nobreakdash}{\nobreak\hbox{-}}
\setlength{\emergencystretch}{2em}
\multlinegap0pt
\usepackage{amssymb}
\usepackage{mathtools}
\usepackage{xcolor}
\usepackage[shortlabels]{enumitem}
\usepackage{algorithm}\usepackage{algpseudocode}
\usepackage{tikz}
\newtheorem{theorem}{Theorem}
\title{LAWS: Learning from Actual Workloads Symbolically\\ \large A Self-Certifying Parametrized Cache Architecture\\ for Neural Inference, Robotics, and Edge Deployment}
\author{Gregory Magarshak}
\date{Source: arXiv:2605.04069v1 (CC BY 4.0)}
\begin{document}
\maketitle

\noindent\emph{Source and licence.} LAWS: Learning from Actual Workloads Symbolically\\ \large A Self-Certifying Parametrized Cache Architecture\\ for Neural Inference, Robotics, and Edge Deployment. Version arXiv:2605.04069v1 (CC BY 4.0), distributed by arXiv under the Creative Commons Attribution 4.0 International licence, http://creativecommons.org/licenses/by/4.0/, as recorded in the arXiv licence metadata for this exact version and shown on the abstract page https://arxiv.org/abs/2605.04069v1. Full licence notice: \texttt{licenses/arxiv-2605.04069-CC-BY-4.0.txt}.\par\smallskip

\noindent\textit{Editorial scope.} Counted unit: the article's theorem thm:abort (source0.tex lines 964--977). The article's complete original proof of it is delivered with it (source0.tex lines 979--1018, one \texttt{proof} environment of 40 lines). No mathematics is written, repaired or re-derived: the statement and the proof are the article's own lines, in the article's own notation, and no retained line is substituted or re-flowed.  Retained uncounted as the source-local context the proof needs: lines 763--779.  Nothing was omitted from any retained span, so the package carries no omission record.  No retained line contains a citation, so the wrapper carries no bibliography and no bibliography entry.  No retained line contains a figure environment, an image-inclusion command or drawing code, so no image is delivered and the package declares no asset dependency.  Editorial formatting only, in addition: an AMS \texttt{amsart} preamble that restates the article's own declarations and macros used by the retained lines, namely the environments \texttt{theorem} on separate counters, together with the standard packages those lines need.  The retained environments print in this document as Theorem 1 in order of appearance, whereas the article prints the same items with the numbering of its own full document.  The complete native source of the article as distributed in the arXiv e-print for this exact version is bundled unchanged as \texttt{sources/199834/source0.tex}.
\begin{theorem}[Jacobian Correction Achieves Second-Order Error]
\label{thm:jacobian}
Let $e$ be a Level-2 expert with $f(\phi) = F_W(n^*) + J_W(n^*) \cdot \phi$
and $\phi(x) = E(x_{\bar{d}+1}) - E(n^*_{\bar{d}+1})$ where $\bar{d}$ is the
length of the longest common prefix of $x$ and $n^*$.  Define
$r = \|E(x_{\bar{d}+1}) - E(n^*_{\bar{d}+1})\|$ as the embedding-space distance
at the divergence point.  Then:
\[
  \|F_W(x) - f(\phi(x))\|
  \;\leq\;
  \frac{1}{2} \|H_W(n^*)\|_{\mathrm{op}} \cdot r^2
  + O(r^3)
\]
where $H_W(n^*)$ is the Hessian of $F_W$ at $n^*$ with respect to the divergence
embedding.  In particular the error is $O(r^2)$ versus $O(r)$ for a Level-1
(constant) expert, giving a strict improvement for $r < 1$.
\end{theorem}

\begin{theorem}[Abort-and-Replan Threshold]
\label{thm:abort}
Let $C_{\mathrm{full}} > 0$ be the cost of a full base-model forward pass,
$C_{\mathrm{hit}} > 0$ the cost of an expert evaluation, and $\lambda \geq 0$
the downstream cost per unit output error.  The optimal validity radius that
minimizes expected total cost is:
\[
  \tau^* = \frac{2(C_{\mathrm{full}} - C_{\mathrm{hit}})}{\lambda \cdot
  \|H_W\|_{\mathrm{op}}^2 \cdot \Lambda(W)^2 \cdot C_E^2},
\]
where $\|H_W\|_{\mathrm{op}}$ is the Hessian operator norm at the signpost.
At $d_{\mathcal{T}}(x, n^*) = \tau^*$, the marginal cost of using the cache
equals the marginal cost of full inference.
\end{theorem}

\begin{proof}
The expected total cost when using a Level-2 expert with routing radius $\tau$ is:
\[
  \mathrm{Cost}(\tau) = \underbrace{C_{\mathrm{hit}} \cdot P(\text{hit})}_{
    \text{cheap path}} +
  \underbrace{C_{\mathrm{full}} \cdot P(\text{miss})}_{
    \text{expensive path}} +
  \underbrace{\lambda \cdot \mathbb{E}[\text{error}^2 \mid \text{hit}]}_{
    \text{error cost}}.
\]

\textbf{Error cost model.}  By Theorem~\ref{thm:jacobian}, for a hit at query $x$
with divergence embedding $r = \|E(x_{\bar{d}+1}) - E(n^*_{\bar{d}+1})\|$,
the squared error is $\leq \frac{1}{4}\|H_W\|_{\mathrm{op}}^2 \cdot r^4$.
We model the expected squared $r$ for queries routed to the expert as increasing
with $\tau$: concretely, for the local approximation $\mathbb{E}[r^2 \mid
\text{hit}] \approx (\Lambda(W) C_E)^2 \cdot \tau^2$, which holds when the
conditional distribution of $r$ given $d_{\mathcal{T}}(x, n^*) \leq \tau$
concentrates at $r \approx \Lambda(W) C_E \cdot \tau$ (the Lipschitz-propagated
trie distance).  Under this approximation:
\[
  \mathbb{E}[\text{error}^2 \mid \text{hit}]
  \;\approx\; \tfrac{1}{4}\|H_W\|_{\mathrm{op}}^2 \cdot (\Lambda(W) C_E)^2 \cdot \tau^2.
\]

\textbf{Miss probability model.}  $P(\text{miss}) = 1 - P_{\mathcal{M}}(\mathcal{B}(n^*, \tau))
\approx 1 - c\tau$ for small $\tau$, where $c = \lim_{\tau\to 0}
P_{\mathcal{M}}(\mathcal{B}(n^*, \tau))/\tau$ is the local $P_{\mathcal{M}}$-density
in trie metric space.

\textbf{Optimization.}  Differentiating $\mathrm{Cost}(\tau)$ and setting to zero:
\[
  -(C_{\mathrm{full}} - C_{\mathrm{hit}}) \cdot c
  + \tfrac{\lambda}{2}\|H_W\|_{\mathrm{op}}^2 \cdot (\Lambda(W) C_E)^2 \cdot \tau
  = 0,
\]
giving $\tau^* = 2(C_{\mathrm{full}} - C_{\mathrm{hit}}) \cdot c /
(\lambda \cdot \|H_W\|_{\mathrm{op}}^2 \Lambda(W)^2 C_E^2)$.  For the
canonical case $c = 1$ this simplifies to the stated formula.
\end{proof}

\end{document}
